\documentclass[11pt, letterpaper]{article}
\usepackage[ruled,vlined]{algorithm2e}
% --- Encoding & Margins ---
\usepackage[utf8]{inputenc}
\usepackage[T1]{fontenc}
\usepackage{geometry}
\usepackage{framed}
\usepackage[backend=biber]{biblatex}
\addbibresource{assignment4.bib}
\geometry{margin=1in}

% --- Fonts ---
\usepackage{mathptmx} % Times New Roman
\usepackage[scaled]{helvet} 
\usepackage{courier}

% --- Graphics & Tables ---
\usepackage{graphicx}
\usepackage{float}
\usepackage{caption}
\usepackage{subcaption}
\usepackage{booktabs}
\usepackage{array}

% --- Math ---
\usepackage{amsmath}
\usepackage{amsfonts}
\usepackage{amssymb}

% --- Code Highlighting ---
\usepackage{listings}
\usepackage{xcolor}

% Professional ICLR-style Syntax Colors
\definecolor{py_comment}{rgb}{0.4, 0.4, 0.4}
\definecolor{py_keyword}{rgb}{0.0, 0.0, 0.7}
\definecolor{py_string}{rgb}{0.0, 0.5, 0.0}
\definecolor{py_num}{rgb}{0.6, 0.1, 0.1}
\definecolor{backcolour}{rgb}{0.98, 0.98, 0.98}
\colorlet{shadecolor}{orange!15}

\lstset{
    backgroundcolor=\color{backcolour},   
    commentstyle=\color{py_comment}\itshape,
    keywordstyle=\color{py_keyword}\bfseries,
    numberstyle=\tiny\color{gray},
    stringstyle=\color{py_string},
    basicstyle=\ttfamily\footnotesize,
    breakatwhitespace=false,         
    breaklines=true,                 
    captionpos=b,                    
    keepspaces=true,                 
    numbers=left,                    
    numbersep=8pt,                  
    showspaces=false,                
    showstringspaces=false,
    showtabs=false,                  
    tabsize=4,
    language=Python,
    frame=single,
    rulecolor=\color{lightgray}
}

% --- Links & Headers ---
\usepackage{hyperref}
\usepackage{fancyhdr}

\pagestyle{fancy}
\fancyhf{}
\rhead{\thepage}
\lhead{CSCI 8820: Assignment 4}
\renewcommand{\headrulewidth}{0.5pt}

% ============================================================
% DOCUMENT START
% ============================================================
\begin{document}

\begin{titlepage}
    \centering
    \vspace*{1.5in}
    \rule{\linewidth}{0.5mm} \\[0.4cm]
    {\huge \bfseries Assignment 4: Multiscale Edge Detection \par}
    \rule{\linewidth}{0.5mm} \\[1.5cm]
    
    \vspace{0.3in}
    {\Large \textbf{Aditya Bhargava} \par}
    {\large UGA ID: 811235738 \par}
    
    \vfill
    {\large CSCI 8820: Computer Vision and Pattern Recognition \par}
    {\large School of Computing, University of Georgia \par}
    {\large March 29, 2026 \par}
\end{titlepage}

\newpage
\tableofcontents
\newpage

\section{Introduction}
This report details the implementation of a multi-scale Laplacian-of-Gaussian (LoG) edge detection algorithm, referred to as Edge Focusing. By tracking the zero-crossings from a coarse scale ($\sigma = 5.0$) down to a fine scale ($\sigma = 1.0$). The performance is verified against a single-scale LoG approach using Pratt's Figure of Merit (PFOM).

\section{Algorithm and Methodology}

\subsection{LoG Kernel}
To ensure the magnitude of the convolution response is comparable across varying scales, the LoG kernel is normalized by a $\sigma^2$ factor. The kernel is evaluated over the strict range $[-4\sigma, +4\sigma]$ to guarantee an odd-sized convolution mask:
\begin{equation}
    \text{LoG}(x, y) = -\frac{1}{\pi \sigma^2} \left[ 1 - \frac{x^2 + y^2}{2\sigma^2} \right] \exp\left(-\frac{x^2 + y^2}{2\sigma^2}\right)
\end{equation}
To maintain computational efficiency at large scales (e.g., $\sigma = 5.0$), the 2D convolution is performed utilizing the Fast Fourier Transform (FFT).

\subsection{Multi-Scale Laplacian of Gaussian Edge Detection Algorithm}
The Edge Focusing algorithm operates on the assumption that significant structural edges exist at coarser scales, whereas noise does not. 
% A relative variance threshold (anchored at 5\% of the maximum LoG response) is applied to isolate legitimate zero-crossings. 

Once zero-crossings are detected at scale $\sigma_i$, a boolean mask is created. When processing the subsequent finer scale ($\sigma_{i-0.5}$), the algorithm strictly constraints the search for zero-crossings to this masked region.


\begin{algorithm}[H]
\DontPrintSemicolon
\SetKwInOut{Input}{input}
\SetKwInOut{Output}{output}

\Input{Normalized Grayscale Image $I$, $\sigma_{start}=5.0$, $\sigma_{end}=1.0$, step $\Delta\sigma=0.5$}
\Output{Final edge map $E$, intermediate maps for all $\sigma$}

Initialize $SearchMask \leftarrow \text{None}$\;
Initialize dictionary $SavedMaps$\;

\For{$\sigma \leftarrow \sigma_{start}$ \KwTo $\sigma_{end}$ \textbf{step} $-\Delta\sigma$}{
    $K \leftarrow \text{GenerateKernel}(\sigma, \text{range}=[-4\sigma, 4\sigma])$\;
    $L \leftarrow \text{Convolve2D\_FFT}(I, K)$\;
    
    $L_{norm} \leftarrow L / \max(|L|)$\;
    $T \leftarrow \text{CalculateThreshold}(I)$\;
    
    $E_{\sigma} \leftarrow \text{FindZeroCrossings}(L_{norm}, T)$\;
    
    \If{$SearchMask \neq \text{None}$}{
        $E_{\sigma} \leftarrow E_{\sigma} \cap SearchMask$
    }
    
    \If{$\sigma \in \{5.0, 4.0, 3.0, 2.0, 1.0\}$}{
        $SavedMaps[\sigma] \leftarrow E_{\sigma}$\;
    }
    
    $SearchMask \leftarrow \text{Dilate}(E_{\sigma}, \text{iterations}=2)$ \tcp*{Update window}
}

\Return $SavedMaps$\;
\caption{Multi-Scale Edge Detection using Laplacian of Gaussian}
\end{algorithm}

\section{Experimental Verification}
To mathematically verify that Edge Focusing is superior to a Single-Scale LoG operator, the final output must be evaluated against a proxy ground truth (OpenCV Canny optimized via Otsu thresholding). \textbf{Pratt's Figure of Merit (PFOM)} is calculated using Euclidean distance transforms. A PFOM score closer to $1.0$ indicates perfect overlap, while lower scores mathematically penalize excessive noise or localization error. 

\section{Experimental Results}

% ---------------------------------------------------------
% TEST 1 RESULTS
% ---------------------------------------------------------
\subsection{Results: test1.img}
The single-scale $\sigma=1.0$ operator captures the edges but remains susceptible to noise, resulting in a PFOM of 0.5905. Conversely, the $\sigma=5.0$ operator suppresses the noise but suffers from severe localization error (PFOM: 0.5507). The multi-scale edge focusing algorithm successfully tracks the edge down to precise locations and achieving a PFOM of 0.8484.

\begin{figure}[H]
    \centering
    \includegraphics[width=\textwidth]{assignment4_ver6/multi_scale/test1_edge_focusing_ver6.pdf}
    \caption{Multi-Scale Edge Focussing Progression for \texttt{test1.img}}
\end{figure}
\newpage

\begin{figure}[H]
    \centering
    \begin{subfigure}{0.19\textwidth}
        \includegraphics[width=\textwidth]{assignment4_ver6/single_scale/test1/test1_single_scale_5.png}
        \caption*{$\sigma=5.0$}
    \end{subfigure}
    \hfill
    \begin{subfigure}{0.19\textwidth}
        \includegraphics[width=\textwidth]{assignment4_ver6/single_scale/test1/test1_single_scale_4.png}
        \caption*{$\sigma=4.0$}
    \end{subfigure}
    \hfill
    \begin{subfigure}{0.19\textwidth}
        \includegraphics[width=\textwidth]{assignment4_ver6/single_scale/test1/test1_single_scale_3.png}
        \caption*{$\sigma=3.0$}
    \end{subfigure}
    \hfill
    \begin{subfigure}{0.19\textwidth}
        \includegraphics[width=\textwidth]{assignment4_ver6/single_scale/test1/test1_single_scale_2.png}
        \caption*{$\sigma=2.0$}
    \end{subfigure}
    \hfill
    \begin{subfigure}{0.19\textwidth}
        \includegraphics[width=\textwidth]{assignment4_ver6/single_scale/test1/test1_single_scale_1.png}
        \caption*{$\sigma=1.0$}
    \end{subfigure}
    \caption{Single-Scale LoG Edge Maps progression for \texttt{test1.img}}
\end{figure}

\begin{figure}[H]
    \centering
    \includegraphics[width=\textwidth]{assignment4_ver6/pfom_opencv/test1_merged_verification.pdf}
    \vspace{0.2cm}
    \includegraphics[width=0.4\textwidth]{assignment4_ver6/pfom_opencv/test1_pfom_table.pdf}
    \caption{PFOM Verification: Single-Scale vs. Multi-Scale for \texttt{test1.img}}
\end{figure}

% ---------------------------------------------------------
% TEST 2 RESULTS
% ---------------------------------------------------------
\subsection{Results: test2.img}

The second image suffers from heavy background shadows. The single-scale $\sigma=1.0$ operator is highly vulnerable to this, yielding the lowest relative PFOM of 0.4855. Meanwhile the multi-scale edge focusing initialized at $\sigma=5.0$ effectively suppresses background noise while accurately localizing the primary edge of the foreground objects, increasing the PFOM score to 0.7884.

\newpage

\begin{figure}[H]
    \centering
    \includegraphics[width=\textwidth]{assignment4_ver6/multi_scale/test2_edge_focusing_ver6.pdf}
    \caption{Multi-Scale Edge Focussing Progression for \texttt{test2.img}}
\end{figure}

\begin{figure}[H]
    \centering
    \begin{subfigure}{0.19\textwidth}
        \includegraphics[width=\textwidth]{assignment4_ver6/single_scale/test2/test2_single_scale_5.png}
        \caption*{$\sigma=5.0$}
    \end{subfigure}
    \hfill
    \begin{subfigure}{0.19\textwidth}
        \includegraphics[width=\textwidth]{assignment4_ver6/single_scale/test2/test2_single_scale_4.png}
        \caption*{$\sigma=4.0$}
    \end{subfigure}
    \hfill
    \begin{subfigure}{0.19\textwidth}
        \includegraphics[width=\textwidth]{assignment4_ver6/single_scale/test2/test2_single_scale_3.png}
        \caption*{$\sigma=3.0$}
    \end{subfigure}
    \hfill
    \begin{subfigure}{0.19\textwidth}
        \includegraphics[width=\textwidth]{assignment4_ver6/single_scale/test2/test2_single_scale_2.png}
        \caption*{$\sigma=2.0$}
    \end{subfigure}
    \hfill
    \begin{subfigure}{0.19\textwidth}
        \includegraphics[width=\textwidth]{assignment4_ver6/single_scale/test2/test2_single_scale_1.png}
        \caption*{$\sigma=1.0$}
    \end{subfigure}
    \caption{Single-Scale LoG Edge Maps progression for \texttt{test2.img}}
\end{figure}

\begin{figure}[H]
    \centering
    \includegraphics[width=0.9\textwidth]{assignment4_ver6/pfom_opencv/test2_merged_verification.pdf}
    \vspace{0.2cm}
    \includegraphics[width=0.4\textwidth]{assignment4_ver6/pfom_opencv/test2_pfom_table.pdf}
    \caption{PFOM Verification: Single-Scale vs. Multi-Scale for \texttt{test2.img}}
\end{figure}

% ---------------------------------------------------------
% TEST 3 RESULTS
% ---------------------------------------------------------
\subsection{Results: test3.img}
The single-scale $\sigma=1.0$ operator detects background texture (PFOM: 0.5198), whereas the single-scale $\sigma=5.0$ operator filters the texture but poorly localizes the sharp edges of the shelves (PFOM: 0.6724). The multi-scale edge focusing algorithm resolves this trade-off resulting in the highest recorded PFOM score of 0.8675.

\begin{figure}[H]
    \centering
    \includegraphics[width=0.9\textwidth]{assignment4_ver6/multi_scale/test3_edge_focusing_ver6.pdf}
    \caption{Multi-Scale Edge Focussing Progression for \texttt{test3.img}}
\end{figure}

\newpage
\begin{figure}[H]
    \centering
    \begin{subfigure}{0.19\textwidth}
        \includegraphics[width=\textwidth]{assignment4_ver6/single_scale/test3/test3_single_scale_5.png}
        \caption*{$\sigma=5.0$}
    \end{subfigure}
    \hfill
    \begin{subfigure}{0.19\textwidth}
        \includegraphics[width=\textwidth]{assignment4_ver6/single_scale/test3/test3_single_scale_4.png}
        \caption*{$\sigma=4.0$}
    \end{subfigure}
    \hfill
    \begin{subfigure}{0.19\textwidth}
        \includegraphics[width=\textwidth]{assignment4_ver6/single_scale/test3/test3_single_scale_3.png}
        \caption*{$\sigma=3.0$}
    \end{subfigure}
    \hfill
    \begin{subfigure}{0.19\textwidth}
        \includegraphics[width=\textwidth]{assignment4_ver6/single_scale/test3/test3_single_scale_2.png}
        \caption*{$\sigma=2.0$}
    \end{subfigure}
    \hfill
    \begin{subfigure}{0.19\textwidth}
        \includegraphics[width=\textwidth]{assignment4_ver6/single_scale/test3/test3_single_scale_1.png}
        \caption*{$\sigma=1.0$}
    \end{subfigure}
    \caption{Single-Scale LoG Edge Maps progression for \texttt{test3.img}}
\end{figure}

\begin{figure}[H]
    \centering
    \includegraphics[width=\textwidth]{assignment4_ver6/pfom_opencv/test3_merged_verification.pdf}
    \vspace{0.2cm}
    \includegraphics[width=0.4\textwidth]{assignment4_ver6/pfom_opencv/test3_pfom_table.pdf}
    \caption{PFOM Verification: Single-Scale vs. Multi-Scale for \texttt{test3.img}}
\end{figure}

\section{Conclusion and Analysis}
Based on the experimental verification, the Multi-Scale LoG operator (Edge Focusing) is superior to the Single-Scale LoG operator. 

As mathematically proven by Pratt's Figure of Merit (PFOM), applying a single-scale LoG at $\sigma=1.0$ results in excellent spatial localization but is susceptible to noise. Conversely, applying a single-scale LoG at $\sigma=5.0$ eliminates noise but causes structural edges to drift apart significantly from their true locations (localization error). 

% By utilizing a coarse-to-fine tracking mask, the edge detection algorithm inherits the noise-immunity of the $\sigma=5.0$ scale, while retaining the precise pixel localization of the $\sigma=1.0$ scale. 

Across all three test images, the multi-scale approach yielded the highest PFOM scores when benchmarked against the proxy ground truth (Canny edge detector).

\begin{shaded}
\textbf{\textit{Note:}} For source code, kindly refer to the appendix on the next page.
\end{shaded}

\begin{shaded}
\textbf{\textit{Note:}} The code uses OpenCV for canny edge detection algorithm, which serves as baseline for comparing the single scale and multiscale edge detection using pratt figure of merit (PFOM).
\end{shaded}

\newpage
\section{Appendix: Source Code}
% Paste your final python code here using the lstlisting environment.
\begin{lstlisting}[language=Python]
import numpy as np
import matplotlib.pyplot as plt
import os
import cv2
from scipy.ndimage import distance_transform_edt

plt.rcParams.update(
    {
        "font.family": "serif",
        "font.size": 10,
        "axes.titlesize": 12,
        "figure.dpi": 300,
    }
)


def load_and_preprocess(path):
    try:
        with open(path, "rb") as f:
            _ = f.read(512)
            img_data = np.fromfile(f, dtype=np.uint8)
        return img_data.reshape((512, 512)).astype(float)
    except Exception as e:
        print(f"Error loading file: {e}")
        exit()


def calc_stats(image):

    mean_val = np.mean(image)
    std_val = np.std(image)

    if std_val > 0:
        skewness = np.mean(((image - mean_val) / std_val) ** 3)
    else:
        skewness = 0

    s_norm = min(1.0, abs(skewness))
    mul = max(0.3, 1.0 - s_norm)

    T_temp = mean_val
    for _ in range(50):
        G1 = image[image > T_temp]
        G2 = image[image <= T_temp]
        mu1 = np.mean(G1) if len(G1) > 0 else 0
        mu2 = np.mean(G2) if len(G2) > 0 else 0

        T_new = (mu1 + mu2) / 2.0
        if abs(T_new - T_temp) < 1.0:
            break
        T_temp = T_new

    return mean_val, std_val, T_temp, mul


def generate_log_kernel(sigma):

    radius = int(np.ceil(4 * sigma))

    x = np.arange(-radius, radius + 1)
    y = np.arange(-radius, radius + 1)
    X, Y = np.meshgrid(x, y)

    arg = -(X**2 + Y**2) / (2 * sigma**2)

    kernel = -(1.0 / (np.pi * sigma**2)) * (1.0 + arg) * np.exp(arg)

    kernel -= np.mean(kernel)
    return kernel


def convolve2d_fft(image, kernel):
    shape = (image.shape[0] + kernel.shape[0] - 1, image.shape[1] + kernel.shape[1] - 1)

    fft_img = np.fft.fft2(image, s=shape)
    fft_ker = np.fft.fft2(kernel, s=shape)

    result = np.real(np.fft.ifft2(fft_img * fft_ker))

    offset_y = kernel.shape[0] // 2
    offset_x = kernel.shape[1] // 2

    return result[
        offset_y : offset_y + image.shape[0], offset_x : offset_x + image.shape[1]
    ]


def dilate_mask(mask, iterations=2):
    dilated = np.copy(mask)
    shifts = [(0, 1), (0, -1), (1, 0), (-1, 0), (1, 1), (1, -1), (-1, 1), (-1, -1)]

    for _ in range(iterations):
        current_iter_mask = np.copy(dilated)
        for dy, dx in shifts:
            current_iter_mask |= np.roll(np.roll(dilated, dy, axis=0), dx, axis=1)
        dilated = current_iter_mask

    return dilated


def zero_crossings(log_img, search_mask=None, thresh_map=0.0):

    right = np.roll(log_img, -1, axis=1)
    bottom = np.roll(log_img, -1, axis=0)
    br = np.roll(np.roll(log_img, -1, axis=0), -1, axis=1)
    bl = np.roll(np.roll(log_img, -1, axis=0), 1, axis=1)

    zc_h = (log_img * right < 0) & (np.abs(log_img - right) > thresh_map)
    zc_v = (log_img * bottom < 0) & (np.abs(log_img - bottom) > thresh_map)
    zc_d1 = (log_img * br < 0) & (np.abs(log_img - br) > thresh_map)
    zc_d2 = (log_img * bl < 0) & (np.abs(log_img - bl) > thresh_map)

    crossings = zc_h | zc_v | zc_d1 | zc_d2

    if search_mask is not None:
        crossings = crossings & search_mask

    crossings[:2, :] = False
    crossings[-2:, :] = False
    crossings[:, :2] = False
    crossings[:, -2:] = False

    return crossings


def generate_opencv_benchmark(image):
    img_uint8 = np.clip(image, 0, 255).astype(np.uint8)
    blurred = cv2.GaussianBlur(img_uint8, (5, 5), 1.0)
    high_thresh, _ = cv2.threshold(blurred, 0, 255, cv2.THRESH_BINARY + cv2.THRESH_OTSU)
    low_thresh = 0.5 * high_thresh
    canny_edges = cv2.Canny(blurred, low_thresh, high_thresh)
    return canny_edges > 0


def pratt_figure_of_merit(detected_edges, true_edges, alpha=1.0 / 9.0):

    N_I = np.sum(true_edges)
    N_A = np.sum(detected_edges)

    if N_A == 0 or N_I == 0:
        return 0.0

    inverse_true = np.logical_not(true_edges)
    distances = distance_transform_edt(inverse_true)

    d_i = distances[detected_edges > 0]

    pfom = np.sum(1.0 / (1.0 + alpha * (d_i**2)))
    pfom /= max(N_I, N_A)

    return pfom


def run_single_scale_log(norm_image, image, sigma, T_iterative, mul):
    p = np.ones_like(image, dtype=float)
    p[image <= T_iterative] = 1.5
    p[image > T_iterative] = 0.8

    kernel = generate_log_kernel(sigma)
    log_img = convolve2d_fft(norm_image, kernel)

    max_log_val = np.max(np.abs(log_img))
    if max_log_val > 0:
        log_img = log_img / max_log_val

    base_thresh = 0.05 * mul
    adaptive_thresh_map = base_thresh * p

    return zero_crossings(
        log_img, search_mask=None, thresh_map=adaptive_thresh_map
    )


def run_edge_focusing(image):
    mean_val, std_val, T_iterative, mul = calc_stats(image)
    norm_image = (image - mean_val) / (std_val + 1e-5)

    p = np.ones_like(image, dtype=float)
    p[image <= T_iterative] = 1.5
    p[image > T_iterative] = 0.8

    sigmas_to_process = np.arange(5.0, 0.5, -0.5)
    sigmas_to_save = [5.0, 4.0, 3.0, 2.0, 1.0]

    saved_edge_maps = {}
    current_search_mask = None

    for sigma in sigmas_to_process:
        kernel = generate_log_kernel(sigma)
        log_img = convolve2d_fft(norm_image, kernel)

        max_log_val = np.max(np.abs(log_img))
        if max_log_val > 0:
            log_img = log_img / max_log_val

        base_thresh = 0.05 * mul
        adaptive_thresh_map = base_thresh * p

        edges = zero_crossings(
            log_img,
            search_mask=current_search_mask,
            thresh_map=adaptive_thresh_map,
        )

        if sigma in sigmas_to_save:
            saved_edge_maps[sigma] = edges

        current_search_mask = dilate_mask(edges, iterations=2)

    return saved_edge_maps


def plot_and_save_results(image, edge_maps, filename, output_dir):
    fig, axes = plt.subplots(2, 3, figsize=(15, 10))
    axes = axes.ravel()

    axes[0].imshow(image, cmap="gray")
    axes[0].set_title(f"Original: {filename}")
    axes[0].axis("off")

    plot_idx = 1
    for sigma in [5.0, 4.0, 3.0, 2.0, 1.0]:
        axes[plot_idx].imshow(edge_maps[sigma], cmap="gray")
        axes[plot_idx].set_title(f"Edge Map (\u03c3 = {sigma})")
        axes[plot_idx].axis("off")
        plot_idx += 1

    plt.tight_layout()
    out_path = os.path.join(output_dir, f"{filename.split('.')[0]}_edge_focusing.pdf")
    plt.savefig(out_path, format="pdf", bbox_inches="tight")
    plt.close()


if __name__ == "__main__":
    base_dir = "assignment4_figures"
    os.makedirs(base_dir, exist_ok=True)
    dir_multi = os.path.join(base_dir, "multi_scale")
    dir_single = os.path.join(base_dir, "single_scale")
    dir_pfom = os.path.join(base_dir, "pfom_opencv")

    os.makedirs(dir_multi, exist_ok=True)
    os.makedirs(dir_single, exist_ok=True)
    os.makedirs(dir_pfom, exist_ok=True)

    files = ["test1.img", "test2.img", "test3.img"]

    for file in files:
        if not os.path.exists(file):
            continue

        img = load_and_preprocess(file)
        base_name = file.split(".")[0]

        mean_val, std_val, T_iterative, dyn_mult = calc_stats(img)
        norm_image = (img - mean_val) / (std_val + 1e-5)

        multi_scale_maps = run_edge_focusing(img)
        plot_and_save_results(img, multi_scale_maps, file, dir_multi)
        final_multi_edges = multi_scale_maps[1.0]

        img_single_dir = os.path.join(dir_single, base_name)
        os.makedirs(img_single_dir, exist_ok=True)

        single_scale_maps = {}
        for s in [5.0, 4.0, 3.0, 2.0, 1.0]:
            ss_edges = run_single_scale_log(norm_image, img, s, T_iterative, dyn_mult)
            single_scale_maps[s] = ss_edges
            plt.imsave(
                os.path.join(img_single_dir, f"{base_name}_single_scale_{int(s)}.png"),
                ss_edges,
                cmap="gray",
            )

        ground_truth = generate_opencv_benchmark(img)
        pfom_single_1 = pratt_figure_of_merit(single_scale_maps[1.0], ground_truth)
        pfom_single_5 = pratt_figure_of_merit(single_scale_maps[5.0], ground_truth)
        pfom_multi = pratt_figure_of_merit(final_multi_edges, ground_truth)

        fig, axes = plt.subplots(1, 4, figsize=(20, 5))

        axes[0].imshow(single_scale_maps[1.0], cmap="gray")
        axes[0].set_title(f"Single-Scale (\u03c3=1.0)\nPFOM: {pfom_single_1:.4f}")
        axes[0].axis("off")

        axes[1].imshow(single_scale_maps[5.0], cmap="gray")
        axes[1].set_title(f"Single-Scale (\u03c3=5.0)\nPFOM: {pfom_single_5:.4f}")
        axes[1].axis("off")

        axes[2].imshow(final_multi_edges, cmap="gray")
        axes[2].set_title(f"Multi-Scale Focusing (\u03c3=1.0)\nPFOM: {pfom_multi:.4f}")
        axes[2].axis("off")

        axes[3].imshow(ground_truth, cmap="gray")
        axes[3].set_title("Ground Truth (OpenCV)")
        axes[3].axis("off")

        plt.tight_layout()
        plt.savefig(
            os.path.join(dir_pfom, f"{base_name}_single_multi_opencv.pdf"),
            format="pdf",
            bbox_inches="tight",
        )
        plt.close()

        fig_tbl, ax_tbl = plt.subplots(figsize=(6, 2))
        ax_tbl.axis("tight")
        ax_tbl.axis("off")

        table_data = [
            ["Method", "PFOM Score"],
            ["Single-Scale (\u03c3=1.0)", f"{pfom_single_1:.4f}"],
            ["Single-Scale (\u03c3=5.0)", f"{pfom_single_5:.4f}"],
            ["Multi-Scale (\u03c3=1.0)", f"{pfom_multi:.4f}"],
        ]

        table = ax_tbl.table(cellText=table_data, loc="center", cellLoc="center")
        table.auto_set_font_size(False)
        table.set_fontsize(12)
        table.scale(1, 1.5)

        plt.savefig(
            os.path.join(dir_pfom, f"{base_name}_pfom_table.pdf"),
            format="pdf",
            bbox_inches="tight",
        )
        plt.close(fig_tbl)

\end{lstlisting}

\end{document}